% Bank soal — bab 4
% Bab ditentukan oleh nama berkas ini.

\begin{exo}[Mengintegralkan bentuk diferensial sepanjang suatu lintasan]{integrale-forme-differentielle}
\keywords{bentuk diferensial, diferensial eksak, kriteria Schwarz, integral garis}

Tinjaulah dua bentuk diferensial
\[
\omega_1 = y\,\mathrm{d}x + x\,\mathrm{d}y,
\qquad
\omega_2 = y\,\mathrm{d}x,
\]
dan tiga lintasan yang menghubungkan titik asal $O(0,0)$ dengan titik $M(1,1)$:
lintasan $\gamma_1$ terdiri atas ruas mendatar $O\to(1,0)$ lalu ruas tegak
$(1,0)\to M$; lintasan $\gamma_2$ terdiri atas ruas tegak $O\to(0,1)$ lalu
ruas mendatar $(0,1)\to M$; dan lintasan $\gamma_3$ adalah ruas garis lurus
yang langsung menghubungkan $O$ dan $M$.

\begin{enumerate}
\item Tunjukkan bahwa $\omega_1$ eksak dengan menentukan fungsi $f$ sehingga $\omega_1=\mathrm{d}f$.
\item Hitung $\int_{\gamma_1}\omega_1$ dan $\int_{\gamma_2}\omega_1$ dengan memparameterkan setiap ruas. Temukan kembali hasilnya tanpa perhitungan.
\item Terapkan kriteria Schwarz pada $\omega_2$. Apa kesimpulannya?
\item Hitung $\int_{\gamma_1}\omega_2$, $\int_{\gamma_2}\omega_2$, dan $\int_{\gamma_3}\omega_2$. Berikan komentar.
\end{enumerate}

\begin{indication}
Pada ruas mendatar $\mathrm{d}y=0$, sedangkan pada ruas tegak
$\mathrm{d}x=0$: setiap integral menyusut menjadi integral sederhana.
\end{indication}

\begin{solution}
\textbf{Pertanyaan 1.} Fungsi $f(x,y)=xy$ memenuhi syarat karena $\partial f/\partial x=y$ dan $\partial f/\partial y=x$.

\textbf{Pertanyaan 2.} Untuk kurva yang diparameterkan oleh
$t\mapsto\bigl(x(t),y(t)\bigr)$, berlaku
\[
\int_\gamma\omega_1
= \int \left[y(t)x'(t)+x(t)y'(t)\right]\,\mathrm{d}t.
\]

Lintasan $\gamma_1$ merupakan gabungan dua ruas. Pada ruas pertama,
$x(t)=t$ dan $y(t)=0$, dengan $t\in[0,1]$; jadi $\mathrm{d}x=\mathrm{d}t$
dan $\mathrm{d}y=0$. Pada ruas kedua, $x(t)=1$ dan $y(t)=t$, masih dengan
$t\in[0,1]$; jadi $\mathrm{d}x=0$ dan $\mathrm{d}y=\mathrm{d}t$. Maka
\[
\begin{aligned}
I_1
= \int_{\gamma_1}\omega_1
&= \int_0^1\left(0\times 1+t\times 0\right)\,\mathrm{d}t
 + \int_0^1\left(t\times 0+1\times 1\right)\,\mathrm{d}t \\
&= 0+\int_0^1\mathrm{d}t = 1.
\end{aligned}
\]

Untuk $\gamma_2$, ruas pertama diparameterkan dengan $x(t)=0$, $y(t)=t$,
lalu ruas kedua dengan $x(t)=t$, $y(t)=1$, dan $t\in[0,1]$ pada keduanya.
Karena itu,
\[
\begin{aligned}
I_2
= \int_{\gamma_2}\omega_1
&= \int_0^1\left(t\times 0+0\times 1\right)\,\mathrm{d}t
 + \int_0^1\left(1\times 1+t\times 0\right)\,\mathrm{d}t \\
&= 0+\int_0^1\mathrm{d}t = 1.
\end{aligned}
\]
Kedua nilai sama. Hal ini dapat diperkirakan tanpa perhitungan: karena
$\omega_1=\mathrm{d}f$ dengan $f(x,y)=xy$, integralnya hanya bergantung pada
titik ujung,
\[
\int_\gamma\omega_1=f(M)-f(O)=f(1,1)-f(0,0)=1.
\]

\textbf{Pertanyaan 3.} Tuliskan $\omega_2=A(x,y)\,\mathrm{d}x+B(x,y)\,\mathrm{d}y$.
Di sini $A(x,y)=y$ dan $B(x,y)=0$. Jika bentuk tersebut eksak, kriteria Schwarz
akan memberikan
\[
\frac{\partial A}{\partial y}=\frac{\partial B}{\partial x}.
\]
Namun $\partial A/\partial y=1$, sedangkan $\partial B/\partial x=0$.
Turunan silang berbeda, sehingga $\omega_2$ tidak eksak. Integral dapat
bergantung pada lintasan antara $O$ dan $M$, sebagaimana dibuktikan berikut.

\textbf{Pertanyaan 4.} Karena $\omega_2=y\,\mathrm{d}x$, hanya perubahan $x$
pada nilai $y$ yang tidak nol yang dapat menyumbang pada integral. Dengan
parameterisasi sebelumnya, untuk $\gamma_1$ diperoleh
\[
\begin{aligned}
J_1
=\int_{\gamma_1}\omega_2
&=\int_0^1 0\times 1\,\mathrm{d}t
  +\int_0^1 t\times 0\,\mathrm{d}t \\
&=0.
\end{aligned}
\]
Pada $\gamma_2$, ruas pertama memenuhi $x(t)=0$, jadi $\mathrm{d}x=0$;
pada ruas kedua, $x(t)=t$, $y(t)=1$, dan $\mathrm{d}x=\mathrm{d}t$. Maka
\[
\begin{aligned}
J_2
=\int_{\gamma_2}\omega_2
&=\int_0^1 t\times 0\,\mathrm{d}t
  +\int_0^1 1\times 1\,\mathrm{d}t \\
&=1.
\end{aligned}
\]
Terakhir, ruas diagonal $\gamma_3$ diparameterkan secara eksplisit oleh
\[
x(t)=t,\qquad y(t)=t,\qquad t\in[0,1],
\]
sehingga $\mathrm{d}x=\mathrm{d}t$ dan $\mathrm{d}y=\mathrm{d}t$. Jadi
\[
J_3
=\int_{\gamma_3}\omega_2
=\int_0^1 y(t)x'(t)\,\mathrm{d}t
=\int_0^1 t\,\mathrm{d}t
=\frac12.
\]
Ketiga lintasan memiliki titik ujung yang sama, tetapi memberikan nilai
berbeda: $J_1=0$, $J_2=1$, dan $J_3=1/2$. Inilah ketergantungan pada lintasan
yang menjadi ciri bentuk diferensial tak eksak.
\end{solution}
\end{exo}

\begin{exo}[Perubahan energi dalam yang sama, tiga cara transfer]{meme-delta-u-trois-transferts}
\keywords{hukum pertama, energi dalam, kalor, kerja mekanik, kerja listrik, kalorimetri, fungsi keadaan}

Suatu cairan, kalorimeternya, dan sebuah hambatan kecil yang tercelup membentuk
sistem tertutup yang secara makroskopik diam. Wadahnya kaku dan kapasitas kalor
total sistem dianggap konstan:
\[
C=2{,}00\ \mathrm{kJ\,K^{-1}}.
\]
Sistem hendak dipindahkan dari keadaan setimbang $A$, dengan
$T_A=20{,}0\,^\circ\mathrm{C}$, ke keadaan setimbang $B$, dengan
$T_B=25{,}0\,^\circ\mathrm{C}$. Diberikan bahwa
\[
U(B)-U(A)=C(T_B-T_A).
\]
Perubahan energi kinetik dan energi potensial makroskopik sama dengan nol.
Semua rugi energi ke lingkungan diabaikan.

Ketiga protokol berikut dilakukan secara terpisah dari keadaan awal $A$ yang sama:
\begin{itemize}
\item[Protokol 1] Sistem dikontakkan secara termal dengan sumber panas tanpa menerima kerja apa pun.
\item[Protokol 2] Dinding diisolasi secara termal. Sebuah pengaduk mengaduk cairan karena jatuhnya massa $m$ sejauh $h=5{,}00$ m. Pengaduk dan cairan akhirnya diam, dan seluruh energi mekanik yang hilang dari massa dipindahkan ke sistem. Gunakan $g=9{,}81\ \mathrm{m\,s^{-2}}$.
\item[Protokol 3] Sistem menerima kalor $Q_3=4{,}00$ kJ dan kekurangan energinya diberikan kepada hambatan secara listrik. Daya listrik yang diterima konstan, $\mathcal P=300$ W.
\end{itemize}

\begin{enumerate}
\item Hitung perubahan energi dalam $\Delta U=U(B)-U(A)$ yang sama bagi ketiga protokol.
\item Untuk masing-masing dari dua protokol pertama, tentukan kalor $Q$ dan kerja $W$ yang diterima sistem. Pada protokol kedua, hitung massa $m$ yang diperlukan.
\item Untuk protokol ketiga, hitung kerja listrik yang diterima dan lama hambatan dialiri listrik.
\item Ketiga proses memiliki keadaan awal dan akhir yang sama. Jelaskan mengapa nilai $Q$ dan $W$ tetap dapat berbeda. Apakah sistem mengandung “kalor” atau “kerja” pada keadaan akhir $B$?
\end{enumerate}

\begin{indication}
Gunakan $\Delta U=Q+W$ dengan konvensi ala bankir. Untuk pengaduk, kerja yang
diterima sama dengan penurunan $mgh$ energi potensial massa. Energi listrik
yang melintasi batas melalui kabel dihitung sebagai kerja listrik.
\end{indication}

\begin{solution}
\textbf{Pertanyaan 1.} Selisih suhu ialah $T_B-T_A=5{,}0$ K. Oleh karena itu,
\[
\Delta U=C(T_B-T_A)
=2{,}00\times5{,}0
=10{,}0\ \mathrm{kJ}.
\]
Nilai ini hanya bergantung pada dua keadaan setimbang $A$ dan $B$.

\textbf{Pertanyaan 2.} Pada protokol pertama tidak ada kerja yang diterima:
$W_1=0$. Hukum pertama memberikan
\[
Q_1=\Delta U=10{,}0\ \mathrm{kJ}.
\]
Kalor positif karena energi memasuki sistem.

Pada protokol kedua, dinding terisolasi secara termal, sehingga $Q_2=0$.
Seluruh perubahan energi dalam berasal dari kerja mekanik pengaduk:
\[
W_2=\Delta U=10{,}0\ \mathrm{kJ}.
\]
Karena $W_2=mgh$, diperoleh
\[
m=\frac{W_2}{gh}
=\frac{10{,}0\times10^3}{9{,}81\times5{,}00}
\simeq 204\ \mathrm{kg}.
\]
Besarnya orde nilai ini mengingatkan bahwa kenaikan suhu yang kecil saja
sudah bersesuaian dengan energi mekanik yang besar.

\textbf{Pertanyaan 3.} Hukum pertama mensyaratkan
\[
W_3=\Delta U-Q_3=10{,}0-4{,}00=6{,}00\ \mathrm{kJ}.
\]
Kerja ini positif: catu listrik memberikan energi kepada sistem. Dengan
$W_3=\mathcal P\,\Delta t$,
\[
\Delta t=\frac{W_3}{\mathcal P}
=\frac{6{,}00\times10^3}{300}
=20{,}0\ \mathrm{s}.
\]

\textbf{Pertanyaan 4.} Ketiga lintasan masing-masing memberikan
\[
(Q_1,W_1)=(10{,}0,0)\ \mathrm{kJ},
\qquad
(Q_2,W_2)=(0,10{,}0)\ \mathrm{kJ},
\]
dan
\[
(Q_3,W_3)=(4{,}00,6{,}00)\ \mathrm{kJ}.
\]
Dalam semua kasus, jumlah $Q+W$ adalah $10{,}0$ kJ dan menghasilkan perubahan
$\Delta U$ yang sama. Energi dalam merupakan fungsi keadaan, sedangkan kalor
dan kerja menggambarkan transfer selama proses. Pada keadaan akhir $B$, sistem
memiliki energi dalam, tetapi tidak mengandung “kalor” ataupun “kerja”.
\end{solution}
\end{exo}

\begin{exo}[Kompresi pada tekanan luar konstan]{compression-pression-exterieure-constante}
\seoready{true}
\keywords{kerja gaya tekanan, tekanan luar konstan, kompresi, ekspansi, ekspansi bebas}

Suatu gas dikompresi dari volume $V_A$ ke volume $V_B<V_A$ di bawah tekanan
luar konstan $P_0$.

\begin{enumerate}
\item Hitung kerja yang diterima gas.
\item Gas kembali dari $V_B$ ke $V_A$ melawan tekanan luar $P_0$ yang sama. Hitung kerja yang diterima dan bandingkan dengan kerja kompresi.
\item Ulangi ekspansi dari $V_B$ ke $V_A$ ketika gas mengembang ke ruang hampa. Tafsirkan ketiga tanda yang diperoleh.
\end{enumerate}

\begin{indication}
Gunakan $W=-\int P_{\mathrm{ext}}\,\mathrm{d}V$. Yang harus diintegralkan ialah
tekanan \emph{luar}, termasuk ketika proses tidak kuasistatik.
\end{indication}

\begin{solution}
\textbf{Pertanyaan 1.} Tekanan luar konstan, sehingga
\[
W_{A\to B}=-P_0(V_B-V_A)=P_0(V_A-V_B)>0.
\]
Gas menerima kerja selama kompresi.

\textbf{Pertanyaan 2.} Untuk proses kembali,
\[
W_{B\to A}=-P_0(V_A-V_B)<0.
\]
Kedua kerja saling berlawanan karena kedua proses berlangsung melawan tekanan
luar konstan yang sama.

\textbf{Pertanyaan 3.} Dalam ruang hampa, $P_{\mathrm{ext}}=0$, sehingga
\[
W_{B\to A}^{\mathrm{vide}}=0.
\]
Jadi ekspansi tidak otomatis berarti kerja negatif: harus benar-benar ada
gaya luar yang didorong.
\end{solution}
\end{exo}
